% ---------------------------------------------------------------------
%  CHƯƠNG 1 — MỞ ĐẦU
% ---------------------------------------------------------------------
\chapter{Mở đầu}
\label{ch:modau}

\begin{muctieu}
Chương này nêu bối cảnh của bài toán điểm danh ở giảng đường, lý do chọn
hướng giải quyết bằng nhận diện khuôn mặt, các mục tiêu cụ thể mà đồ án cam
kết đạt được, phạm vi thực hiện và bố cục của báo cáo.
\end{muctieu}

\section{Lý do chọn đề tài}

Một tiết học ở giảng đường thường bắt đầu bằng việc gọi tên hoặc chuyền tay
tờ danh sách ký. Với lớp học phần vài chục sinh viên, việc này lấy mất
khoảng năm tới mười phút đầu giờ, tức gần một phần mười thời lượng tiết học.
Tính trên toàn bộ một học kỳ, phần thời gian đó là đáng kể.

Vấn đề thứ hai nghiêm trọng hơn thời gian: điểm danh thủ công không xác thực
được danh tính. Một sinh viên có thể ký thay bạn, hoặc dạ thay khi giảng viên
gọi tên. Bản thân tờ danh sách ký cũng không để lại bằng chứng nào kiểm chứng
được về sau, nên khi có khiếu nại về kết quả chuyên cần, nhà trường gần như
không có căn cứ để đối chiếu.

Vấn đề thứ ba nằm ở khâu tổng hợp. Dữ liệu điểm danh nằm rải rác trên giấy
hoặc trong các tệp bảng tính riêng của từng giảng viên, nên thông thường chỉ
đến cuối kỳ nhà trường mới biết sinh viên nào vắng quá tỉ lệ cho phép. Khi
đó sinh viên đã mất quyền dự thi và mọi nhắc nhở đều là quá muộn. Phụ huynh
lại càng không có kênh nào để theo dõi tình hình học tập của con trong kỳ.

Nhận diện khuôn mặt là hướng giải quyết hợp lý cho ba vấn đề trên, vì khuôn
mặt là đặc điểm sinh trắc luôn mang theo người, không cần phát thêm thẻ,
không cần chạm vào thiết bị, và một khung hình chụp cả lớp có thể nhận ra
nhiều người cùng lúc. Quan trọng hơn, các mô hình nhận diện hiện đại \cite{Wang2021Survey} cho
phép thêm một người mới chỉ bằng một tấm ảnh, không phải huấn luyện lại, nên
hệ thống mở rộng được theo từng kỳ tuyển sinh. Đây cũng là bài toán cho phép
em vận dụng trọn vẹn các kiến thức về học sâu, thị giác máy tính và tìm kiếm
trên không gian vector đã học trong học phần.

\begin{luuy}
Đổi lại, dữ liệu khuôn mặt là dữ liệu sinh trắc, thuộc nhóm dữ liệu cá nhân
nhạy cảm \cite{Luat912025, ND3562025}. Một hệ thống điểm danh bằng khuôn mặt chỉ có ý nghĩa thực tiễn nếu
đi kèm cơ chế xin phép, phân quyền truy cập và quyền xoá dữ liệu. Ràng buộc
này được coi là yêu cầu bắt buộc của đồ án, chứ không phải phần làm thêm.
\end{luuy}

\section{Mục tiêu của đồ án}
\label{sec:muctieu}

Đồ án đặt ra năm mục tiêu cụ thể, mỗi mục tiêu đều kiểm chứng được và sẽ
được đối chiếu lại trong Chương~\ref{ch:ketqua}:

\begin{enumerate}[label=\textbf{MT\arabic*.}, leftmargin=3.0em, itemsep=3pt]
  \item \textbf{Xây dựng được quy trình nhận diện khuôn mặt hoàn chỉnh} từ
        đăng ký tới điểm danh, chạy được trên máy chủ chỉ có CPU, thời gian
        xử lý một khung hình đủ nhanh để dùng trực tiếp trong lớp.
  \item \textbf{Chống được ba kiểu gian lận thường gặp}: điểm danh bằng ảnh
        in hoặc ảnh chụp lại màn hình, nhận nhầm danh tính do một khung hình
        xấu, và điểm danh từ xa khi không có mặt tại lớp.
  \item \textbf{Số hoá trọn vẹn nghiệp vụ quản lý chuyên cần}: hồ sơ sinh
        viên, lớp học phần, lịch học, phiên điểm danh, sửa điểm danh có lưu
        vết, cảnh báo vắng, thông báo cho phụ huynh và xuất báo cáo, kèm nhật
        ký hệ thống ghi lại mọi thao tác quản lý.
  \item \textbf{Phân quyền đúng mô hình tổ chức của trường} cho quản trị
        viên, Phòng Đào tạo, Phòng Công tác sinh viên, văn phòng Khoa, lãnh
        đạo Khoa, ban ngành, giảng viên, sinh viên và phụ huynh, trong đó cán
        bộ Khoa chỉ thấy dữ liệu thuộc Khoa mình và người nhập liệu tách khỏi
        người phê duyệt.
  \item \textbf{Cung cấp một ứng dụng dùng chung cho Web, Android và iOS}
        từ một mã nguồn duy nhất, kèm bộ kiểm thử tự động chạy trên cơ sở dữ
        liệu thật để bảo vệ các luồng nghiệp vụ quan trọng.
\end{enumerate}

\section{Phạm vi và đối tượng}

\textbf{Về bài toán trí tuệ nhân tạo}, đồ án sử dụng mô hình nhận diện khuôn
mặt đã được tiền huấn luyện thay vì huấn luyện một mô hình mới từ đầu. Lý do
được phân tích ở Mục~\ref{sec:vosaokhonghuanluyen}. Phần đóng góp của đồ án
nằm ở quy trình bao quanh mô hình: kiểm tra chất lượng ảnh, chống giả mạo,
xác nhận đa khung hình, so khớp theo phạm vi lớp và đối chiếu vị trí.

\textbf{Về nghiệp vụ}, đồ án số hoá quy trình điểm danh và quản lý chuyên
cần. Hệ thống không quản lý điểm học phần, học phí, đăng ký tín chỉ hay
nhân sự, và không kết nối với phần mềm quản lý đào tạo sẵn có của trường.

\textbf{Về người dùng}, hệ thống phục vụ \FsoVaiTro{} vai trò được liệt kê ở
Mục~\ref{sec:phanquyen}. Ứng dụng web hướng tới cán bộ và giảng viên, còn
ứng dụng di động hướng tới sinh viên và phụ huynh.

\textbf{Những gì đồ án không làm}: không nhận diện khi sinh viên đeo khẩu
trang che quá nửa khuôn mặt; không xử lý được trường hợp anh chị em sinh đôi
cùng trứng bằng riêng khuôn mặt; không triển khai trên hạ tầng thật của nhà
trường; và chưa kiểm định mô hình trên tập dữ liệu khuôn mặt người Việt quy
mô lớn.

\section{Phương pháp thực hiện}

Đồ án được thực hiện theo hướng phát triển lặp \cite{Larman2003Iterative}. Mỗi vòng lặp gồm bốn bước:
chọn một luồng nghiệp vụ, cài đặt luồng đó ở cả phía dịch vụ lẫn ứng dụng,
viết kiểm thử tự động cho luồng, rồi chạy thử với dữ liệu mẫu để tìm lỗi
thực tế. Cách làm này giữ cho hệ thống luôn ở trạng thái chạy được, thay vì
dồn toàn bộ phần tích hợp vào cuối kỳ.

Toàn bộ mã nguồn được quản lý bằng \code{git}. Tại thời điểm viết báo cáo,
kho mã nguồn có \FsoCommit{} commit, bản được mô tả trong báo cáo này là
commit \code{\FcommitCuoi} ngày \FngayCommit.

\begin{ghinho}
Mọi con số về quy mô hệ thống trong báo cáo này, từ số bảng dữ liệu tới số
ca kiểm thử, đều được sinh tự động từ chính kho mã nguồn bằng công cụ mô tả
ở mục C của phần Phụ lục, nên không có số nào là ước lượng bằng tay.
\end{ghinho}

\section{Bố cục báo cáo}

Phần còn lại của báo cáo được tổ chức như sau. Chương~\ref{ch:coso} trình
bày cơ sở lý thuyết về nhận diện khuôn mặt, chống giả mạo, tìm kiếm trên
không gian vector và các công nghệ được sử dụng. Chương~\ref{ch:thietke}
phân tích yêu cầu và thiết kế kiến trúc, cơ sở dữ liệu, giao diện lập trình
cùng mô hình phân quyền. Chương~\ref{ch:caidat} mô tả quá trình cài đặt các
thành phần cốt lõi và giao diện người dùng. Chương~\ref{ch:ketqua} trình bày
kết quả đạt được, kết quả kiểm thử, đánh giá so với mục tiêu ban đầu, rồi
tổng kết, nêu bài học rút ra và hướng phát triển.
